\chapter{Savage Acts and Subjective Expected-Utility Representation}
\label{sec:p17-savage-acts}

The finite monetary-bet theorem in
\cref{sec:p17-finite-subjective-probability} averages payments
under a strong vector-additivity assumption. The present problem
is different: derive both an event probability and a possibly
nonlinear consequence utility from preference over acts.
No probability or utility is supplied at the outset.
The construction follows the atomless branch of Abdellaoui and
Wakker's 2020 primary reconstruction, including its probability
consistency argument \cite{AbdellaouiWakker2020SavageSA}.
The original Savage postulates and the later arbitrary-act
formulation are distinguished using Frahm and Hartmann
\cite{FrahmHartmann2025SavageSA}.

\section{States, events, acts and postulates}
\label{sec:p17-sa-domain}

\begin{hbdefinition}{def:p17-sa-acts}{Act domains and conditional preference}
Let \(S\ne\varnothing\) be a state set, \(X\ne\varnothing\)
a consequence set, and \(\mathcal E\) an algebra of subsets
of \(S\). A simple act is an \(\mathcal E\)-measurable
finite-valued map \(f:S\to X\). All such acts are available.
Neither \(S\) nor \(X\) is assumed finite.
Write \(f_Ag\) for the act equal to \(f\) on \(A\) and
\(g\) on \(A^c\), and identify a consequence with its constant
act. A weak order \(\succeq\) has strict part \(\succ\)
and indifference \(\sim\).

Conditional preference \(f\succeq_Ag\) means
\(f_Ah\succeq g_Ah\) for a common background \(h\).
Under P2 below it is independent of that background.
An event is null if every pair of acts is conditionally
indifferent on it. Empty events are null.
For the arbitrary-act formulation later, take
\(\mathcal E=2^S\) and \(\mathcal F=X^S\), not merely
the simple acts.
\end{hbdefinition}

\begin{hbdefinition}{def:p17-sa-postulates}{Savage's P1--P7 with their domains}
The first five conditions are:
\begin{enumerate}
\item P1: preference is complete and transitive.
\item P2, sure thing:
\(f_Ah\succeq g_Ah\) iff \(f_Ak\succeq g_Ak\)
for every available \(f,g,h,k,A\).
\item P3: on every nonnull event \(A\),
\(x\succeq y\) iff \(x\succeq_Ay\).
\item P4: for any \(x\succ y\) and \(v\succ w\),
\(x_Ay\succeq x_By\) iff \(v_Aw\succeq v_Bw\).
\item P5: some constant consequences satisfy \(x\succ y\).
\end{enumerate}
P6, small-event continuity, says that if \(f\succ g\)
and \(x\in X\), a finite event partition \(\{A_i\}\) of
\(S\) exists with \(x_{A_i}f\succ g\) and
\(f\succ x_{A_i}g\) for every cell.
The replacement is a constant consequence, not an arbitrary
new act. Those are distinct formulations.

On \(\mathcal F=X^S\), P7 consists of both implications
\[
 [\,f(s)\succeq_Ag\ \text{for every }s\in A\,]
       \ \Longrightarrow\ f\succeq_Ag,
\]
\[
 [\,f\succeq_Ag(s)\ \text{for every }s\in A\,]
       \ \Longrightarrow\ f\succeq_Ag.
\]
The comparisons involving \(f(s)\) or \(g(s)\) are comparisons
with constant acts, not statewise comparisons between numerical
utilities already assigned.
\end{hbdefinition}

These are source-specific postulates, not consequences of the
word rational. The 2020 paper uses a constant-background
version of P2 on simple acts. Successive replacement on its
finite background partition gives the form above; its converse
is immediate. P7 concerns the extension beyond that domain.
The 1972 revision's bounded-utility qualification must not be
silently read into the 1954 manifestation.

\begin{hbdefinition}{def:p17-sa-construction-axioms}{The 2020 atomless construction assumptions}
Fix P1--P5 on all simple acts. In place of P6 impose:
\begin{enumerate}
\item Archimedeanity (P6.1): no infinite sequence of disjoint,
equally likely nonnull events exists.
\item Event solvability (P6.2): if \(x\succ y\) and
\(y_Bf\prec g\prec x_Bf\), there is \(C\subseteq B\)
in \(\mathcal E\) with
\(g\sim x_Cy_{B\setminus C}f\).
\end{enumerate}
Define \(A\succeq_eB\) by \(x_Ay\succeq x_By\)
for one \(x\succ y\). P4 makes the stakes irrelevant.
For this branch also assume no qualitative atoms:
every \(A\succ_e\varnothing\) contains
\(\varnothing\prec_e B\prec_e A\).
\end{hbdefinition}

Solvability is an additional richness assumption, not an
effective procedure for eliciting its event. The paper's
atomic branch has different restrictions and is not substituted
for this construction. Its Appendix B derives the new assumptions
from Savage's P6 on a sigma-algebra; the algebra-only theorem
and the original sigma-algebra presentation remain distinct.

\section{Constructing probability from almost uniform partitions}
\label{sec:p17-sa-probability}

\begin{hbdefinition}{def:p17-sa-cup}{Close-to-uniform partitions and cuts}
A CUP is a partition
\(\{A_1,\ldots,A_n,R\}\) of \(S\), \(n\geq1\), with
all \(A_i\sim_e A_1\succ_e\varnothing\) and
\(R\prec_e A_1\); the remainder may be empty.
Write \([j]_A=\bigcup_{i=1}^j A_i\), \([0]_A=\varnothing\).
For an event \(E\), let \(k\) be the largest
\(0\leq j\leq n\) with \([j]_A\preceq_e E\).
Use the conservative numerical interval
\[
 I_A(E)=\left[\frac{k-1}{n+1},\frac{k+2}{n}\right].
\]
This weak lower-cut convention differs at equality by at
most one unit from the source's strict lower cut; the
conservative-margin method is the same and its limit is unchanged.
These numbers are constructed from comparisons, not assumed
probabilities.
\end{hbdefinition}

\begin{hbtheorem}{thm:p17-sa-cup-consistency}{proved=annex:p17-sa-cup-consistency}{CUP construction and consistency of probability cuts}
Under the stated atomless construction assumptions, CUPs exist
with arbitrarily many equal cells. For every fixed coarse CUP,
all sufficiently refined CUP intervals lie strictly inside its
conservative interval. The intersection over all CUPs is a
single point \(P(E)\).
The resulting \(P:\mathcal E\to[0,1]\) is normalized and
finitely additive, represents \(\succeq_e\), and is the
unique finitely additive probability representing that event
order.
\end{hbtheorem}

The critical consistency comparison is explicit in the annex.
If a fine CUP has \(N\) equal cells, \(i\) fit into a coarse
equal cell, and \(\ell\) fit into \(E\), then
\[
 ni\leq N<(n+1)(i+1),\qquad
 ki\leq\ell\leq(k+1)(i+1).
\]
Hence the fine lower and upper bounds tend respectively
past the coarse conservative bounds as \(i\to\infty\).
One cannot merely assert that separate probability estimates
have a common limit. This supplies the source's central
Lemma 10 argument in algebraic form
\cite[pp.\ 8--11, 15--16]{AbdellaouiWakker2020SavageSA}.

\section{Constructing utility and the simple-act representation}
\label{sec:p17-sa-utility}

\begin{hbtheorem}{thm:p17-sa-utility}{proved=annex:p17-sa-utility}{Conditional calibration and utility existence}
Choose consequences \(M\succ m\).
For \(m\preceq x\preceq M\), solvability gives an event
\(S^x\) with \(x\sim M_{S^x}m\).
Set \(U(m)=0\), \(U(M)=1\), \(U(x)=P(S^x)\).
This is independent of the chosen calibrating event.
For every nonnull \(A\), a conditional calibrating
\(A^x\subseteq A\) satisfies
\[
 x\sim_A M_{A^x}m,\qquad P(A^x)=U(x)P(A).
\]
Consequently every simple act
\((E_1:x_1,\ldots,E_r:x_r)\) between those consequences
is indifferent to a two-consequence gamble with winning
probability \(\sum_jP(E_j)U(x_j)\).

Enlarging the consequence anchors and keeping the fixed
normalization at \(m,M\) yields one
\(U:X\to\mathbb R\) on the entire consequence set.
For all simple \(f,g\),
\[
 f\succeq g\quad\Longleftrightarrow\quad
 \sum_jP(E_j)U(x_j)\geq
 \sum_jP(E_j)U(y_j),
\]
where \(f=x_j\) and \(g=y_j\) on a common finite event
partition \(\{E_j\}\).
The utility is unique up to a positive affine transformation.
It need not be globally bounded on this simple-act domain.
\end{hbtheorem}

The conditional identity is not supplied as a state-independent
utility assumption. It is proved from stake-independent event
comparisons and calibration. The annex constructs a dominated
additive set function \(T_x(A)=P(A^x)\), then proves
\(T_x=U(x)P\) using CUPs. This is the source's probability/
utility exchange-rate argument, not monetary vector additivity.
There may be infinitely many consequences across the domain
although each one act takes finitely many
\cite[pp.\ 11--14]{AbdellaouiWakker2020SavageSA}.

\section{Arbitrary acts, bounded utility and the source boundary}
\label{sec:p17-sa-general-acts}

\begin{hbdefinition}{def:p17-sa-integral}{Finite-additive expectation for bounded utility acts}
On \(\mathcal E=2^S\), define the integral of a simple real
function by \(I_P(\sum_jc_j\mathbf1_{E_j})=\sum_jc_jP(E_j)\).
For any bounded real function \(v\), approximate it uniformly
by finite-valued functions and take the limit of their integrals.
For a globally bounded \(u:X\to\mathbb R\), every
arbitrary act \(f\in X^S\) therefore has the finite value
\(I_P(u\circ f)\). Countable additivity is not presumed.
\end{hbdefinition}

\begin{hbtheorem}{thm:p17-sa-general-converse}{proved=annex:p17-sa-general-converse}{Integral construction and the arbitrary-act converse}
The bounded integral exists, is independent of approximants,
is linear and monotone, and has
\(|I_P(v)-I_P(w)|\leq\lVert v-w\rVert_\infty\).
If \(P\) is convex-ranged and \(u\) is bounded and nonconstant,
comparison by \(I_P(u\circ f)\) on all \(X^S\) satisfies
P1--P7, including null events and small-event continuity.
Here convex-ranged means that every event \(A\) contains
a subevent of every probability between zero and \(P(A)\).
\end{hbtheorem}

\begin{hbtheorem}{thm:p17-sa-general-source}{cited=FrahmHartmann2025SavageSA}{The sourced arbitrary-act existence and uniqueness}
On all acts \(X^S\) and all events \(2^S\), P1--P7 imply
a convex-ranged finitely additive probability \(P\) and
a nonconstant globally bounded utility \(u\) representing
preference by \(I_P(u\circ f)\).
Conversely that representation satisfies the postulates.
The probability is unique and utility is unique up to
positive affine change.
\end{hbtheorem}

The cited forward direction includes the extension to nonsimple
acts and necessity of globally bounded utility. It is not
proved merely by the simple-act construction above.
Frahm and Hartmann distinguish this from P1--P6 on simple
acts and discuss the boundedness correction recognized in
Savage's 1972 revision
\cite[Sec.\ 2]{FrahmHartmann2025SavageSA}.
Nonatomicity alone is not convex-rangedness for a general
finitely additive measure. Countable additivity, if desired,
requires an additional continuity condition on decreasing
event sequences, not just P6.

\begin{hbtheorem}{thm:p17-sa-identification}{proved=annex:p17-sa-identification}{What identifies probability and utility}
With all acts available and a convex-ranged probability,
two state-independent SEU representations of the same
preference have identical probabilities and positive-affine
related utilities. A positive probability atom contradicts
P6 when utility is nonconstant.
If state-dependent consequence scales are allowed instead,
that decomposition need not be unique.
\end{hbtheorem}

\begin{hbexample}{def:p17-sa-failure}{Richness and identification failures}
With two consequences and two states, weights \((.7,.3)\)
and \((.8,.2)\) order all four binary acts alike; event
rankings alone do not reveal cardinal probabilities in
that restricted domain. This is a failure of the richness
required for the general construction, not its counterexample.
More generally a nonnull atom cannot be made small by a
finite partition, as the P6 proof shows.

For a rich event space take \(P(E)=1/2\) and bounded
state-independent \(u\). Let
\(Q(A)=1.5P(A\cap E)+.5P(A\cap E^c)\).
Assign utility \(u/1.5\) on \(E\) and \(u/.5\) on \(E^c\).
The expected values of every act are unchanged, yet
\(Q(E)=3/4\). This alternative uses state-dependent scales;
it is not an alternative state-independent representation,
nor a violation of P4 by the original preferences.
\end{hbexample}

\begin{proofboundary}{def:p17-sa-boundaries}{Distinct presentations}
The proved constructive representation here is the atomless,
simple-act branch of the 2020 primary, on general state and
consequence spaces. Its finite-additive integral and converse
also treat arbitrary bounded utility acts.
The source's atomic branch, the original route through lottery
mixtures, and the full forward arbitrary-act extension retain
their separate stated statuses. Finite states, finite-valued
acts, and globally bounded consequence utility are three
different restrictions.
\end{proofboundary}
